% ------------------------------------------------------------

% LaTeX Template für die DHBW zum Schnellstart!
% Original: https://github.wdf.sap.corp/vtgermany/LaTeX-Template-DHBW
% ------------------------------------------------------------
% ---- Präambel mit Angaben zum Dokument
\documentclass[
	fontsize=12pt,           % Leitlinien sprechen von Schriftgröße 12.
	paper=A4,
	twoside=false,
	listof=totoc,            % Tabellen- und Abbildungsverzeichnis ins Inhaltsverzeichnis
	bibliography=totoc,      % Literaturverzeichnis ins Inhaltsverzeichnis aufnehmen
	titlepage,               % Titlepage-Umgebung anstatt \maketitle
	headsepline,             % horizontale Linie unter Kolumnentitel
	abstract,              % Überschrift einschalten, Abstract muss in {abstract}-Umgebung stehen
]{scrreprt}                  % Verwendung von KOMA-Report
\usepackage[utf8]{inputenc}  % UTF8 Encoding einschalten
\usepackage[ngerman]{babel}  % Neue deutsche Rechtschreibung
\usepackage[T1]{fontenc}     % Ausgabe von westeuropäischen Zeichen (auch Umlaute)
\usepackage{microtype}       % Trennung von Wörtern wird besser umgesetzt
\usepackage{lmodern}         % Nicht-gerasterte Schriftarten (bei MikTeX erforderlich)
\usepackage{graphicx}        % Einbinden von Grafiken erlauben
\usepackage{wrapfig}         % Grafiken fließend im Text
\usepackage{setspace}        % Zeilenabstand \singlespacing, \onehalfspaceing, \doublespacing
\usepackage[
	%showframe,                % Ränder anzeigen lassen
	left=2cm, right=2cm,
	top=1.1cm,  bottom=1.1cm,
	includeheadfoot, headsep=1.25cm, footskip=1.25cm
]{geometry}                      % Seitenlayout einstellen
%\usepackage[left=2.5cm, right=2.5cm, head=1.25cm, bottom=1.25cm, foot=1.25cm, includefoot]{geometry}

\usepackage{scrlayer-scrpage}    % Gestaltung von Fuß- und Kopfzeilen
\usepackage{acronym}             % Abkürzungen, Abkürzungsverzeichnis
\usepackage{titletoc}            % Anpassungen am Inhaltsverzeichnis
\contentsmargin{0.75cm}          % Abstand im Inhaltsverzeichnis zw. Punkt und Seitenzahl
\usepackage[                     % Klickbare Links (enth. auch "nameref", "url" Package)
  hidelinks,                     % Blende die "URL Boxen" aus.
  breaklinks=true                % Breche zu lange URLs am Zeilenende um
]{hyperref}
\usepackage[hypcap=true]{caption}% Anker Anpassung für Referenzen

\usepackage{lineno}

\usepackage{textcomp}

\usepackage{chngcntr}
\counterwithout{figure}{chapter}

\renewcommand{\thefigure}{\arabic{figure}}
\renewcommand{\thetable}{\arabic{table}}

\urlstyle{same}                  % Aktuelle Schrift auch für URLs
% Anpassung von autoref für Gleichungen (ergänzt runde Klammern) und Algorithm.
% Anstatt "Listing" kann auch z.B. "Code-Ausschnitt" verwendet werden. Dies sollte
% jedoch synchron gehalten werden mit \lstlistingname (siehe weiter unten).
\addto\extrasngerman{%
	\def\equationautorefname~#1\null{Gleichung~(#1)\null}
	\def\lstnumberautorefname{Zeile}
	\def\lstlistingautorefname{Listing}
	\def\algorithmautorefname{Algorithmus}
	% Damit einheitlich "Abschnitt 1.2[.3]" verwendet wird und nicht "Unterabschnitt 1.2.3"
	% \def\subsectionautorefname{Abschnitt}
}

% ---- Abstand verkleinern von der Überschrift 
\renewcommand*{\chapterheadstartvskip}{\vspace*{.5\baselineskip}}

% Hierdurch werden Schusterjungen und Hurenkinder vermieden, d.h. einzelne Wörter
% auf der nächsten Seite oder in einer einzigen Zeile.
% LaTeX kann diese dennoch erzeugen, falls das Layout ansonsten nicht umsetzbar ist.
% Diese Werte sind aber gute Startwerte.
\widowpenalty10000
\clubpenalty10000

% ---- Für das Quellenverzeichnis
\usepackage[
	backend = biber,                % Verweis auf biber
	language = auto,
	style = numeric,                % Nummerierung der Quellen mit Zahlen
	citestyle=authoryear,
	sorting = none,                 % none = Sortierung nach der Erscheinung im Dokument
	sortcites = true,               % Sortiert die Quellen innerhalb eines cite-Befehls
	block = space,                  % Extra Leerzeichen zwischen Blocks
	hyperref = true,                % Links sind klickbar auch in der Quelle
	backref = true,                % Referenz, auf den Text an die zitierte Stelle
	bibencoding = auto,
	giveninits = true,              % Vornamen werden abgekürzt
	doi=false,                      % DOI nicht anzeigen
	isbn=false,                     % ISBN nicht anzeigen
    alldates=short                  % Datum immer als DD.MM.YYYY anzeigen
]{biblatex}
\addbibresource{Inhalt/literatur.bib}
\setcounter{biburlnumpenalty}{3000}     % Umbruchgrenze für Zahlen
\setcounter{biburlucpenalty}{6000}      % Umbruchgrenze für Großbuchstaben
\setcounter{biburllcpenalty}{9000}      % Umbruchgrenze für Kleinbuchstaben
\DeclareNameAlias{default}{family-given}  % Nachname vor dem Vornamen
\AtBeginBibliography{\renewcommand{\multinamedelim}{\addslash\space
}\renewcommand{\finalnamedelim}{\multinamedelim}}  % Schrägstrich zwischen den Autorennamen
\DefineBibliographyStrings{german}{
  urlseen = {Einsichtnahme:},                      % Ändern des Titels von "besucht am"
}
\usepackage[babel,german=quotes]{csquotes}         % Deutsche Anführungszeichen + Zitate


\usepackage{xcolor}
\usepackage{blindtext}
\usepackage{soul}

% ---- Für Mathevorlage
\usepackage{amsmath}    % Erweiterung vom Mathe-Satz
\usepackage{amssymb}    % Lädt amsfonts und weitere Symbole
\usepackage{MnSymbol}   % Für Symbole, die in amssymb nicht enthalten sind.


% ---- Für Quellcodevorlage
\usepackage{scrhack}                    % Hack zur Verw. von listings in KOMA-Script
\usepackage{listings}                   % Darstellung von Quellcode
\usepackage{xcolor}                     % Einfache Verwendung von Farben
% -- Eigene Farben für den Quellcode
\definecolor{JavaLila}{rgb}{0.4,0.1,0.4}
\definecolor{JavaGruen}{rgb}{0.3,0.5,0.4}
\definecolor{JavaBlau}{rgb}{0.0,0.0,1.0}
\definecolor{ABAPKeywordsBlue}{HTML}{6000ff}
\definecolor{ABAPCommentGrey}{HTML}{808080}
\definecolor{ABAPStringGreen}{HTML}{4da619}
\definecolor{PyKeywordsBlue}{HTML}{0000AC}
\definecolor{PyCommentGrey}{HTML}{808080}
\definecolor{PyStringGreen}{HTML}{008080}
% -- Farben für ABAP CDS
\definecolor{CDSString}{HTML}{FF8C00}
\definecolor{CDSKeywords}{HTML}{6000ff}
\definecolor{CDSAnnotation}{HTML}{00BFFF}
\definecolor{CDSComment}{HTML}{808080}
\definecolor{CDSFunc}{HTML}{FF0000}

% -- Default Listing-Styles

\lstset{
	% Das Paket "listings" kann kein UTF-8. Deswegen werden hier 
	% die häufigsten Zeichen definiert (ä,ö,ü,...)
	literate=%
		{á}{{\'a}}1 {é}{{\'e}}1 {í}{{\'i}}1 {ó}{{\'o}}1 {ú}{{\'u}}1
		{Á}{{\'A}}1 {É}{{\'E}}1 {Í}{{\'I}}1 {Ó}{{\'O}}1 {Ú}{{\'U}}1
		{à}{{\`a}}1 {è}{{\`e}}1 {ì}{{\`i}}1 {ò}{{\`o}}1 {ù}{{\`u}}1
		{À}{{\`A}}1 {È}{{\'E}}1 {Ì}{{\`I}}1 {Ò}{{\`O}}1 {Ù}{{\`U}}1
		{ä}{{\"a}}1 {ë}{{\"e}}1 {ï}{{\"i}}1 {ö}{{\"o}}1 {ü}{{\"u}}1
		{Ä}{{\"A}}1 {Ë}{{\"E}}1 {Ï}{{\"I}}1 {Ö}{{\"O}}1 {Ü}{{\"U}}1
		{â}{{\^a}}1 {ê}{{\^e}}1 {î}{{\^i}}1 {ô}{{\^o}}1 {û}{{\^u}}1
		{Â}{{\^A}}1 {Ê}{{\^E}}1 {Î}{{\^I}}1 {Ô}{{\^O}}1 {Û}{{\^U}}1
		{œ}{{\oe}}1 {Œ}{{\OE}}1 {æ}{{\ae}}1 {Æ}{{\AE}}1 {ß}{{\ss}}1
		{ű}{{\H{u}}}1 {Ű}{{\H{U}}}1 {ő}{{\H{o}}}1 {Ő}{{\H{O}}}1
		{ç}{{\c c}}1 {Ç}{{\c C}}1 {ø}{{\o}}1 {å}{{\r a}}1 {Å}{{\r A}}1
		{€}{{\euro}}1 {£}{{\pounds}}1 {«}{{\guillemotleft}}1
		{»}{{\guillemotright}}1 {ñ}{{\~n}}1 {Ñ}{{\~N}}1 {¿}{{?`}}1,
	breaklines=true,        % Breche lange Zeilen um 
	breakatwhitespace=true, % Wenn möglich, bei Leerzeichen umbrechen
	% Symbol für Zeilenumbruch einfügen
	prebreak=\raisebox{0ex}[0ex][0ex]{\ensuremath{\rhookswarrow}},
	postbreak=\raisebox{0ex}[0ex][0ex]{\ensuremath{\rcurvearrowse\space}},
	tabsize=4,                                 % Setze die Breite eines Tabs
	basicstyle=\ttfamily\small,                % Grundsätzlicher Schriftstyle
	columns=fixed,                             % Besseres Schriftbild
	numbers=left,                              % Nummerierung der Zeilen
	%frame=single,                             % Umrandung des Codes
	showstringspaces=false,                    % Keine Leerzeichen hervorheben
	keywordstyle=\color{blue},
	ndkeywordstyle=\bfseries\color{darkgray},
	identifierstyle=\color{black},
	commentstyle=\itshape\color{JavaGruen},   % Kommentare in eigener Farbe
	stringstyle=\color{JavaBlau},             % Strings in eigener Farbe,
	captionpos=b,                             % Bild*unter*schrift
	xleftmargin=5.0ex
}

% ---- Eigener JAVA-Style für den Quellcode
\renewcommand{\ttdefault}{pcr}               % Schriftart, welche auch fett beinhaltet
\lstdefinestyle{EigenerJavaStyle}{
	language=Java,                             % Syntax Highlighting für Java
	%frame=single,                             % Umrandung des Codes
	keywordstyle=\bfseries\color{JavaLila},    % Keywords in eigener Farbe und fett
	commentstyle=\itshape\color{JavaGruen},    % Kommentare in eigener Farbe und italic
	stringstyle=\color{JavaBlau}               % Strings in eigener Farbe
}

% ---- Eigener ABAP-Style für den Quellcode
\renewcommand{\ttdefault}{pcr}
\lstdefinestyle{EigenerABAPStyle}{
	language=[R/3 6.10]ABAP,
	morestring=[b]\|,                          % Für Pipe-Strings
	morestring=[b]\`,                          % für Backtick-Strings
	keywordstyle=\bfseries\color{ABAPKeywordsBlue},
	commentstyle=\itshape\color{ABAPCommentGrey},
	stringstyle=\color{ABAPStringGreen},
	tabsize=2,
	morekeywords={
		types,
		@data,
		as,
		lower,
		start,
		selection,
		order,
		by,
		inner,
		join,
		key,
		end,
		cast
	}
}

% ---- Eigener Python-Style für den Quellcode
\renewcommand{\ttdefault}{pcr}
\lstdefinestyle{EigenerPythonStyle}{
	language=Python,
	columns=flexible,
	keywordstyle=\bfseries\color{PyKeywordsBlue},
	commentstyle=\itshape\color{PyCommentGrey},
	stringstyle=\color{PyStringGreen}
}

%----- ABAP-CDS-View language
\lstdefinelanguage{ABAPCDS}{
	sensitive=false,
	%Keywords
	morekeywords={define,
		view,
		as,
		select,
		from,
		inner,
		join,
		on,
		key,
		case,
		when,
		then,
		else,
		end,
		true,
		false,
		cast,
		where,
		and,
		distinct,
		group,
		by,
		having,
		min,
		sum,
		max,
		count,
		avg
	},
	%Methoden
	morekeywords=[2]{
		div,
		currency\_conversion,
		dats\_days\_between,
		concat\_with\_space,
		dats\_add_days,
		dats\_is\_valid,
		dats\_add\_months,
		unit\_conversion,
		division,
		mod,
		abs,
		floor,
		ceil,
		round,
		concat,
		replace,
		substring,
		left,
		right,
		length
	},
	morecomment=[s][\color{CDSAnnotation}]{@}{:},
	morecomment=[l][\itshape\color{CDSComment}]{//},
	morecomment=[s][\itshape\color{CDSComment}]{/*}{*/},
	morestring=[b][\color{CDSString}]',
	keywordstyle=\bfseries\color{CDSKeywords},
	keywordstyle=[2]\color{CDSFunc}
}

  % Weitere Details sind ausgelagert

\usepackage{algorithm}                  % Für Algorithmen-Umgebung (ähnlich wie lstlistings Umgebung)
\usepackage{algpseudocode}              % Für Pseudocode. Füge "[noend]" hinzu, wenn du kein "endif",
                                        % etc. haben willst.

\makeatletter                           % Sorgt dafür, dass man @ in Namen verwenden kann.
                                        % Ansonsten gibt es in der nächsten Zeile einen Compilefehler.
\renewcommand{\ALG@name}{Algorithmus}   % Umbenennen von "Algorithm" im Header der Listings.
\makeatother                            % Zeichen wieder zurücksetzen
\renewcommand{\lstlistingname}{Listing} % Erlaubt das Umbenennen von "Listing" in anderen Titel.


% ---- Tabellen
\usepackage{booktabs}  % Für schönere Tabellen. Enthält neue Befehle wie \midrule
\usepackage{multirow}  % Mehrzeilige Tabellen
\usepackage{siunitx}   % Für SI Einheiten und das Ausrichten Nachkommastellen
\sisetup{locale=DE, range-phrase={~bis~}, output-decimal-marker={,}} % Damit ein Komma und kein Punkt verwendet wird.
\usepackage{xfrac} % Für siunitx Option "fraction-function=\sfrac"

% ---- Für Definitionsboxen in der Einleitung
\usepackage{amsthm}                     % Liefert die Grundlagen für Theoreme
\usepackage[framemethod=tikz]{mdframed} % Boxen für die Umrandung
% ---- Definition für Highlight Boxen

% ---- Grundsätzliche Definition zum Style
\newtheoremstyle{defi}
  {\topsep}         % Abstand oben
  {\topsep}         % Abstand unten
  {\normalfont}     % Schrift des Bodys
  {0pt}             % Einschub der ersten Zeile
  {\bfseries}       % Darstellung von der Schrift in der Überschrift
  {:}               % Trennzeichen zwischen Überschrift und Body
  {.5em}            % Abstand nach dem Trennzeichen zum Body Text
  {\thmname{#3}}    % Name in eckigen Klammern
\theoremstyle{defi}

% ------ Definition zum Strich vor eines Texts
\newmdtheoremenv[
  hidealllines = true,       % Rahmen komplett ausblenden
  leftline = true,           % Linie links einschalten
  innertopmargin = 0pt,      % Abstand oben
  innerbottommargin = 4pt,   % Abstand unten
  innerrightmargin = 0pt,    % Abstand rechts
  linewidth = 3pt,           % Linienbreite
  linecolor = gray!40,       % Linienfarbe
]{defStrich}{Definition}     % Name der des formats "defStrich"

% ------ Definition zum Eck-Kasten um einen Text
\newmdtheoremenv[
  hidealllines = true,
  innertopmargin = 6pt,
  linecolor = gray!40,
  singleextra={              % Eck-Markierungen für die Definition
    \draw[line width=3pt,gray!50,line cap=rect] (O|-P) -- +(1cm,0pt);
    \draw[line width=3pt,gray!50,line cap=rect] (O|-P) -- +(0pt,-1cm);
    \draw[line width=3pt,gray!50,line cap=rect] (O-|P) -- +(-1cm,0pt);
    \draw[line width=3pt,gray!50,line cap=rect] (O-|P) -- +(0pt,1cm);
  }
]{defEckKasten}{Definition}  % Name der des formats "defEckKasten"  % Weitere Details sind ausgelagert

% ---- Für Todo Notes
\usepackage{todonotes}
\setlength {\marginparwidth }{2cm}      % Abstand für Todo Notizen

\usepackage[official]{eurosym}

\usepackage{pdfpages}

% ---- Elektronische Version oder Gedruckte Version?
% ---- Unterschied: Die elektronische Version enthält keinen Platzhalter für die Unterschrift
\usepackage{ifthen}
\usepackage{color}
\newboolean{e-Abgabe}
\setboolean{e-Abgabe}{false}    % false=gedruckte Fassung

% ---- Persönlichen Daten:
\newcommand{\titel}{Betriebsplankosten- und -erfolgsrechnung in Industriebetrieben}
\newcommand{\titelheader}{Betriebsplankosten- und -erfolgsrechnung}
\newcommand{\arbeit}{wissenschaftliche Ausarbeitung}
\newcommand{\studiengang}{Wirtschaftsinformatik}
\newcommand{\studienjahr}{2023}
\newcommand{\autor}{Tom Wolfrum}
\newcommand{\autorReverse}{Wolfrum, Tom}
\newcommand{\verfassungsort}{Karlsruhe}
\newcommand{\matrikelnr}{4000776}
\newcommand{\kurs}{WWI22B5}
% \newcommand{\bearbeitungsmonat}{Januar 2018}
\newcommand{\abgabe}{12. November 2023}
\newcommand{\bearbeitungszeitraum}{16.10.2023 - 12.11.2023}
\newcommand{\firmaName}{SAP SE}
\newcommand{\firmaStrasse}{Dietmar-Hopp-Allee 16}
\newcommand{\firmaPlz}{69190 Walldorf, Deutschland}
\newcommand{\betreuerFirma}{Steven Rösinger}
\newcommand{\betreuerDhbw}{Harald Haake}

% ---- Metainformation für das PDF Dokument
\hypersetup{
	pdftitle    = {\titel},
	pdfsubject  = {\arbeit},
	pdfauthor   = {\autor},
	%pdfkeywords = {Keywords angeben},
	pdfcreator  = {LaTeX},
	%pdfproducer = {in der Regel pdfTeX}
}

% ---- Definition der Kopf- und Fußzeilen
\clearpairofpagestyles                          % Löschen von LaTeX Standard
\automark[section]{chapter}                     % Füllen von section und chapter
\renewcommand*{\chaptermarkformat}{}            % Entfernt die Kapitelnummer
\renewcommand*{\sectionmarkformat}{}            % Entfernt die Sectionnummer
% Angaben [für "plain"]{für "scrheadings"}
\ihead[]{\small\titelheader}                          % Kopfzeile links
\chead[]{}                                      % Kopfzeile mitte
\ohead[]{\small\rightmark}                            % Kopfzeile rechts
\ifoot[]{}                                      % Fußzeile links
\cfoot{}                     % Fußzeile mitte
\ofoot*{\sffamily\pagemark}                                      % Fußzeile rechts
\KOMAoptions{
   headsepline = 0.2pt,                         % Liniendicke Kopfzeile
   footsepline = false                          % Liniendicke Fußzeile
}

% ---- Hilfreiches
\newcommand{\zB}{z.\,B. }   % "z.B." mit kleinem Leeraum dazwischen (ohne wäre nicht korrekt)
\newcommand{\dash}{d.\,h. }

\newcommand{\code}[1]{\texttt{#1}} % Ist einfacher zu schreiben als ständig \texttt und erlaubt
                                   % Änderungen im Nachhinein, wenn man z.B. Inline-Code anders stylen möchte.

% ---- Silbentrennung (falls LaTeX defaults falsch / nicht gewünscht sind)
\hyphenation{HANA}         % anstatt HA-NA
\hyphenation{Graph-Script} % anstatt GraphS-cript

% ---- Beginn des Dokuments

\begin{document}
\setlength{\parindent}{0pt}              % Keine Paragraphen Einrückung.
                                         % Dafür haben wir den Abstand zwischen den Paragraphen.
\setcounter{secnumdepth}{2}              % Nummerierungstiefe fürs Inhaltsverzeichnis
\setcounter{tocdepth}{2}                 % Tiefe des Inhaltsverzeichnisses. Ggf. so anpassen,
                                         % dass das Verzeichnis auf eine Seite passt.
\sffamily                                % Serifenlose Schrift verwenden.

% ------ Vorspann
% ------ Titelseite
\singlespacing
\thispagestyle{empty}
\begin{titlepage}
\enlargethispage{4cm}

\begin{figure}           % Logo vom Ausbildungsbetrieb und der DHBW
	% \vspace*{-5mm} % Sollte dein Titel zu lang werden, kannst du mit diesem "Hack" 
	%                  den Inhalt der Seite nach oben schieben.
	\begin{minipage}{0.49\textwidth}
		\flushleft
		% \includegraphics[height=2.5cm]{Bilder/Logos/Logo_SAP.pdf} 
	\end{minipage}
	\hfill
	\begin{minipage}{0.49\textwidth}
		\flushright
		\includegraphics[height=2.5cm]{Bilder/Logos/Logo_DHBW.pdf} 
	\end{minipage}
\end{figure} 
\vspace*{2.5cm}

\begin{center}
	\huge{\textbf{\titel}}\\[1.5cm]
	\Large{\textbf{\arbeit}}\\[0.5cm]
	\normalsize{im Rahmen der Portfolioprüfung in der Vorlesung\\[1ex] \textbf{Geld, Währung und Wirtschaftspolitik}}\\[0.5cm]
	\Large{des Studienganges \studiengang}\\[1ex]
	\normalsize{an der Dualen Hochschule Baden-Württemberg Karlsruhe}\\[1cm]
	\normalsize{von}\\[1ex] \Large{\textbf{\autor}} \\[1cm]

	% Hinweis: Manche Dozenten möchten einen Hinweis auf den Sperrvermerk auf der Titelseite.

	% Sperrvermerkt ein-/auskommentieren:
	% \large{{\color{red}- Sperrvermerk -}}\\[1cm]


\end{center}

\begin{center}
	\vfill
	\begin{tabular}{ll}
		\textbf{Abgabedatum:}                     & \abgabe \\[0.2cm]
		\textbf{Bearbeitungszeitraum:}            & \bearbeitungszeitraum \\[0.2cm]
		\textbf{Kurs:}            				  & \kurs \\[0.2cm]
		\textbf{Ausbildungsfirma:}                & \firmaName \\
		                                 		  & \firmaStrasse \\
		                                          & \firmaPlz \\[0.2cm]
		% \textbf{Betreuer der Ausbildungsfirma:}   & \betreuerFirma \\[0.2cm]
		\textbf{Gutachter der Dualen Hochschule:} & \betreuerDhbw \\[2cm]
	\end{tabular} 
\end{center}
\end{titlepage}
  % Titelseite
\newcounter{savepage}
\pagenumbering{Roman}                    % Römische Seitenzahlen
\onehalfspacing

% ------ Erklärung, Sperrvermerk, Abstact
% \chapter*{Sperrvermerk}
Die nachfolgende Arbeit enthält vertrauliche Daten der:
\begin{quote}
	\firmaName \\
	\firmaStrasse \\
	\firmaPlz
\end{quote}

\vspace{0.5cm}

Der Inhalt dieser Arbeit darf weder als Ganzes noch in Auszügen Personen ausserhalb des Prüfungs- und Evaluationsverfahrens zugänglich gemacht werden, sofern keine anders lautende Genehmigung des Dualen Partners vorliegt.
\chapter*{Selbstständigkeitserklärung}

Ich versichere hiermit, dass ich die vorliegende \arbeit{} mit dem Thema:
\begin{quote}
	\textit{\titel}
\end{quote}
selbstständig verfasst und keine anderen als die angegebenen Quellen und Hilfsmittel benutzt habe.

\vspace{0.25cm}


\vspace{1cm}

\verfassungsort, den \today \\[0.5cm]
\ifthenelse{\boolean{e-Abgabe}}
	{\underline{Gez. \autor}}
	{\makebox[6cm]{\hrulefill}}\\ 
\autorReverse

\chapter*{Geschlechtsneutrale Formulierung}

 

In dieser Arbeit wird aus Gründen der besseren Lesbarkeit das generische Maskulinum verwendet.

Weibliche und anderweitige Geschlechteridentitäten werden dabei ausdrücklich mitgemeint, soweit es für die Aussage erforderlich ist.

%\renewcommand{\abstractname}{Abstract} % Veränderter Name für das Abstract
\begin{abstract}
\begin{addmargin}[1.5cm]{1.5cm}        % Erhöhte Ränder, für Abstract Look
\thispagestyle{plain}                  % Seitenzahl auf der Abstract Seite

\begin{center}
\small\textit{- English -}             % Angabe der Sprache für das Abstract
\end{center}

\vspace{0.25cm}

This is the starting point of the Abstract. For the final bachelor thesis, there must be an abstract included in your document. So, start now writing it in German and English. The abstract is a short summary with around 200 to 250 words.

\vspace{0.25cm}

Try to include in this abstract the main question of your work, the methods you used or the main results of your work.


\end{addmargin}
\end{abstract}
%\renewcommand{\abstractname}{Abstract} % Veränderter Name für das Abstract
\begin{abstract}
\begin{addmargin}[1.5cm]{1.5cm}        % Erhöhte Ränder, für Abstract Look
\thispagestyle{plain}                  % Seitenzahl auf der Abstract Seite

\begin{center}
\small\textit{- Deutsch -}             % Angabe der Sprache für das Abstract
\end{center}

\vspace{0.25cm}

Dies ist der Beginn des Abstracts. Für die finale Bachelorarbeit musst du ein Abstract in deinem Dokument mit einbauen. So, schreibe es am besten jetzt in Deutsch und Englisch. Das Abstract ist eine kurze Zusammenfassung mit ca. 200 bis 250 Wörtern.

\vspace{0.25cm}

Versuche in das Abstract folgende Punkte aufzunehmen: Fragestellung der Arbeit, methodische Vorgehensweise oder die Hauptergebnisse deiner Arbeit.


\end{addmargin}
\end{abstract}

% ------ Inhaltsverzeichnis
\singlespacing

\tableofcontents

% ------ Verzeichnisse
\renewcommand*{\chapterpagestyle}{plain}
\pagestyle{plain}
%\chapter*{Formelverzeichnis}
\addcontentsline{toc}{chapter}{Formelverzeichnis} % Hinzufügen zum Inhaltsverzeichnis 

% Definition des neuen Befehls für das Einfügen der Abkürzung der Einheit
\newcommand{\acrounit}[1]{
  \acroextra{\makebox[18mm][l]{\si[per-mode=fraction,fraction-function=\sfrac]{#1}}}
}
\begin{acronym}[dmin] % längstes Kürzel wird verw. für den Abstand zw. Kürzel u. Text

	% Alphabetisch selbst sortieren - nicht verwendete Formeln rausnehmen!
	% Allgemein: \acro{KÜRZEL}[ABKÜRZUNG]{\acrounit{SI-EINHEIT}BESCHREIBUNG}

	% \acro{A}[\ensuremath{A}]{\acrounit{mm^2}Fläche}	
	% \acro{D}[\ensuremath{D}]{\acrounit{mm}Werkstückdurchmesser}	
	% \acro{dmin}[\ensuremath{d\textsubscript{min}}]{\acrounit{mm}kleinster Schaftdurchmesser}	
	% \acro{L1}[\ensuremath{L\textsubscript{1}}]{\acrounit{mm}Länge des Werkstückes Nr. 1}	
	% \acro{Fwinkel}[]{\acrounit{Grad}Freiwinkel}	
	% \acro{Kwinkel}[]{\acrounit{Grad}Keilwinkel}

\end{acronym}

\chapter*{Abkürzungsverzeichnis}
\addcontentsline{toc}{chapter}{Abkürzungsverzeichnis} % Hinzufügen zum Inhaltsverzeichnis 

\begin{acronym}[WYSISWG] % längstes Kürzel wird verw. für den Abstand zw. Kürzel u. Text

	% Alphabetisch selbst sortieren - nicht verwendete Kürzel rausnehmen!
	
	% Bsp.:
	\acro{T10/B50}{gewichtetes Verhältnis des Durchschnittseinkommens der reichsten 10\% zum Durchschnittseinkommen der ärmsten 50\%}

\end{acronym}
% \listoffigures                          % Erzeugen des Abbildungsverzeichnisses 
\listoftables                           % Erzeugen des Tabellenverzeichnisses
\renewcommand{\lstlistlistingname}{Quellcodeverzeichnis}
%\lstlistoflistings                      % Erzeugen des Listenverzeichnisses
\setcounter{savepage}{\value{page}}


% ------ Inhalt der Arbeit
\cleardoublepage
\pagenumbering{arabic}                  % Arabische Seitenzahlen für den Hauptteil
\setlength{\parskip}{0.5\baselineskip}  % Abstand zwischen Absätzen
\rmfamily
\renewcommand*{\chapterpagestyle}{scrheadings}
\pagestyle{scrheadings}
\singlespacing
%\include{Inhalt/04_Inhalt/einleitung}
%\include{Inhalt/04_Inhalt/formatText}
%\chapter{Technische Grundlagen}
Alle Abbildungen und Tabellen sind laut Richtlinien fortlaufend mit Nummern zu versehen. Diese Aufgabe übernimmt \LaTeX{} für dich. Jedoch solltest du den Text unter dem Bild (Legende) auf das jeweilige Bild anpassen. Hast du das Bild irgendwo entnommen, dann muss der Quellenverweis auch direkt in der Legende mit eingebaut sein.

Im Text selbst solltest du auf die Abbildung verweisen. Neben dem reinen verweisen, solltest du dich auch damit auseinandersetzen. Beschreibe die Grafik, hebe die Relevanz einzelner Teile hervor oder nenne andere wichtige Informationen im Text davor oder danach.

\section{Tabellen}
Die Legende steht bei den Tabellen darüber, während diese bei anderen Grafiken darunter ihren Platz einnimmt.

\begin{table}[ht]
	\centering
	\caption{Eine dreispaltige Tabelle}
	\begin{tabular}{lll}
		\toprule
		\textbf{linke Spalte} & \textbf{mittlere Spalte} & \textbf{rechte Spalte} \\
		\midrule
		A & B & C \\
		! & 2 & 3 \\
		a & b & c \\
		i & ii & iii \\
		\bottomrule
	\end{tabular}
	\label{tabelle:Einfache3Spalten}
\end{table}

Bitte beachte auch, dass \LaTeX{} deine Tabellen an eine andere Position verschiebt, wenn diese dort besser aussehen. Vermeide deshalb Textpassagen wie beispielsweise \enquote{in der nachfolgenden Tabelle/Grafik}, denn es könnte sein, dass die Tabelle vor den Text rutscht. Nutze deshalb immer Verweise wie: \enquote{in der Tabelle \ref{tabelle:Zahlenausrichtung} ist \dots }!

Die Tabelle \ref{tabelle:Zahlenausrichtung} nutzt mehrere Packages. Mithilfe des Befehls \texttt{\textbackslash midrule} aus dem Package \enquote{booktabs} erstellt man eine horizontale Linie, die die vertikalen unterbricht. Diese sorgt für ein schöneres Schriftbild. Möchtest du Zahlen auflisten, so kannst du diese nach dem Dezimalkomma ausrichten. Dies geht mit dem Package \enquote{siunitx}.

\begin{table}[ht]
	\centering
	\caption{Zahlenausrichtung}
	\begin{tabular}{c|c|r|r|S[table-format=3.7]}
		% Anstatt \hline können \toprule, \bottomrule und \midrule verwendet werden.
		\toprule
		% Hinweis: \multicolumn wird verwendet, um die eigentliche Anordnung zu umgehen.
		Nr.  & Datum & Euro & USD  & \multicolumn{1}{c@{}}{Zahlen} \\ 
		\midrule 
		1  & 01.06.2017 & $1,00$€  & $1,13\$$  & 11,158   \\
		2  & 02.06.2017 & $2,00$€  & $2,26\$$  & 2,18     \\
		3  & 03.06.2017 & $3,00$€  & $3,39\$$  & 9,15568  \\
		4  & 04.06.2017 & $4,00$€  & $4,52\$$  & 5,868668 \\
		5  & 05.06.2017 & $5,00$€  & $5,65\$$  & 1,4      \\
		\midrule
		6  & 06.06.2017 & $6,00$€  & $6,78\$$  & 6,58     \\
		7  & 07.06.2017 & $7,00$€  & $7,91\$$  & 7,998    \\
		8  & 08.06.2017 & $8,00$€  & $9,04\$$  & 4,358    \\
		9  & 09.06.2017 & $9,00$€  & $10,17\$$ & 3,5458   \\
		10 & 10.06.2017 & $10,00$€ & $11,30\$$ & 302,8    \\
		\bottomrule
	\end{tabular}
	\label{tabelle:Zahlenausrichtung}
\end{table}

Die Tabelle \ref{tabelle:Zellenverbindung} zeigt die Möglichkeit, einzelne Zellen miteinander zu verbinden. Mit dem Befehl \texttt{\textbackslash multicolumn \{Anzahl Spalten\}\{Ausrichtung\}\{Inhalt\}} kannst du mehrere Spalten verbinden. Dafür müssen die überschriebenen Spalten entfernt werden, es darf also keinen \enquote{\&}-Trennzeichen dafür geben. Beim Verbinden mehrerer Reihen wird der Befehl \texttt{\textbackslash multirow\{Anzahl Reihen\}\{*\}\{Inhalt\}} verwendet. Hierbei darauf achten, dass die Zellen, die zusammengefasst werden sollen, keinen Inhalt haben, aber vorhanden sind, also einen \enquote{\&}-Trennzeichen besitzen.

\begin{table}[ht]
	\centering
  \caption{Verbinden von Zellen}
	\begin{tabular}{c|c|r}
		\toprule
		Text 1               &      Mittiger Text 2      &   Text 3 \\ \midrule
		\multicolumn{2}{l}{Linksbündig}                  &  1,00 \texteuro \\ \midrule
		   2                 &                           &  2,00 €  \\ \midrule
		   3                 &        03.06.2017         &  3,00 € \\ \midrule
		   4                 &        04.06.2017         &  4,00 € \\ \midrule
		                    \multicolumn{3}{r}{Rechtsbündiger Text} \\ \midrule
		   6                 &        06.06.2017         &  6,00 € \\ \midrule
		   7                 & \multirow{2}{*}{Zusammen} &  7,00 € \\ \cmidrule{1-1} \cmidrule{3-3}
		   8                 &                           &  8,00 € \\ \midrule
		   9                 &        09.06.2017         &  9,00 € \\ \midrule
		  10                 &        10.06.2017         & 10,00 € \\ \bottomrule
	\end{tabular}
	\label{tabelle:Zellenverbindung}
\end{table}

\pagebreak
\section{Grafiken}

\begin{wrapfigure}{l}{0.30\textwidth}
	\centering
	\vspace{-20pt} % Manchmal möchte man den oberen Abstand selbst anpassen
	\includegraphics[width=0.25\textwidth]{Bilder/desktop-screen.pdf}
	\vspace{-10pt}
	% Das folgende ist ein Trick, um "Abbilgung x.y" in eine
	% eigene Zeile zu packen. Der Text zwischen [ und ] steht
	% im Abbildungsverzeichnis. Der Text darunter wird
	% tatsächlich angezeigt.
	\caption[Wrap Figure mit einer PDF als Grafik]{\unskip}
	Wrap Figure mit einer PDF als Grafik
	\label{fig:wrap-Referenz-auf-Bild}
\end{wrapfigure}
Eine gute Projektarbeit kommt nicht ohne einige Abbildungen in Form von Skizzen, Diagrammen oder  ähnlichem aus.
Am besten ist es, wenn du diese Grafiken selbst als SVG Dateien erstellst und diese in Form eines PDFs einbindest, somit ist die Grafik auch beim Drucken scharf.

Gerade solche Kleinigkeiten zeugen von Professionalität und können auch nochmals einige Punkte für eine gute Note raus holen. Die \texttt{figure} Umgebung eignet sich für das Einbinden von Grafiken, die die volle Seitenbreite ausnutzen. Sie erscheinen nicht im Fließtext.

Dagegen gibt es die \texttt{wrapfigure} Umgebung, die das Einbinden von Grafiken im Fließtext erlaubt.
Diese Umgebung ist allerdings ab und zu problematisch zu benutzen.
Tipp: Die \texttt{wrapfigure} Umgebung vor dem Absatz platzieren und darauf achten, dass sie nicht in der Nähe eines Seitenumbruchs ist.
Dann erscheint sie rechts oder links des Paragraphen. Mit \texttt{vspace} kann noch der Abstand nach oben und unten angepasst werden, um leeren Platz zu vermeiden.
Siehe auch \url{https://tex.stackexchange.com/questions/56176/handling-of-wrapfig-pictures-in-latex}.

\begin{figure}[ht]
	\centering
	\includegraphics[width=0.50\textwidth]{Bilder/desktop-screen.pdf} 
	\caption{Desktop Ansicht mit einem Symbol}
	\label{fig:Referenz-auf-Bild}
\end{figure}

Das kostenlose Tool \href{https://inkscape.org/}{\textbf{Inkscape}} hilft dir beim erzeugen dieser SVG Grafiken. Solltest du noch nie damit gearbeitet haben, dann schau dir am besten einige kurze Tutorials an. Im Github Repository unter dem Reiter Wiki findest du eine Kurzanleitung, wie du ein PDF in Inkscape generieren kannst. Solltest du weitere gute kostenlose Software für die Bildbearbeitung in SVG kennen, dann kannst du diese gerne uns per E-Mail mitteilen. Danke!

\section{Diagramme, Mockups, Software Design}

Um Software Designs mit einzubinden eignet sich das Online Tool \textbf{\href{https://www.draw.io/}{DRAW.IO}} sehr gut. Mit diesem Tool lassen sich Diagramme, Mockups, Flow Charts, technische Zeichnungen, Wireframe Skizzen sowie Software Designs zeichnen, welche du ohne Verpixelung im PDF-Format wieder, wie in Abbildung \ref{fig:Diagramm-DrawIo} gezeigt, einbinden kannst.

\begin{figure}[ht]
	\centering
	\includegraphics[width=0.8\textwidth]{Bilder/Diagramm.pdf} 
	\caption{Diagramm eines Software Designs}
	\label{fig:Diagramm-DrawIo}
\end{figure}

Ein weiteres kostenloses Tool für die Darstellung von Klassendiagrammen oder auch Sequenzdiagrammen ist \textbf{\href{http://staruml.io/}{StarUML}}. Dieses generiert auch SVG Dateien, welche du mithilfe von Inkscape zu PDF Dateien konvertieren kannst.

%\chapter{Mathematik}

\section{Text}
Hier steht ein beispielhafter Text bei dem nun auf eine sehr bekannte und durchaus vertraute Formel im Text direkt mit $c = \sqrt{a^2 + b^2}$ eingegangen wird. Dabei können auch Winkel wie: $\alpha, \beta, \gamma$ gerne verwendet werden. Weiterhin werden hier nur Formeln im Bereich der $\mathbb{N}$ dargestellt, gerne können diese aber durch Formeln aus diesem Bereich der $\mathbb{R}$ ergänzt werden.

\section{Formeln}
Beachte bei Formeln keine konkreten Werte anzugeben, sondern die Formel stets nur wie in der Literatur nur als Größengleichungen anzugeben.
\begin{align}
	\sum_{n=0}^{\infty}x=b+n\\
	\frac{b*x}{c} = y
\end{align}

Trotz unterschiedlicher Länge kann man die Gleichheitszeichen auf der gleichen Höhe anbringen wie in \autoref{eq:Gleichung1} und \autoref{eq:Gleichung2} dargestellt, zusätzlich kann man diese auch mit Informationen versehen wie in Formel \autoref{eq:Gleichung3} zu sehen.

\begin{align}
	\label{eq:Gleichung1} a + b &= c\\
	\label{eq:Gleichung2} 5c + 3f &= 4h\\ 
	\label{eq:Gleichung3} \overbrace{5y}^{y = 0} + \underbrace{42x}_{x = 1} &= b
\end{align}

\section{Arrays und Matrizen}

\begin{center}
	\(
	\begin{array}{lc|r}
		a&b&c\\
		\hline
		x&y&z\\
		c&a&b
	\end{array}
	\)
\end{center}

\begin{equation}
	\begin{array}{lcl}
		z & = & a \\
		a + b & = & c \\
		f(x,y,z) & = & x + y + z
	\end{array}
\end{equation}

Hier noch ein paar Matrizen Beispiele in \LaTeX{}.

\begin{align}
	\begin{pmatrix}
		a_{11}	& \dots   & a_{1n}\\
		\vdots	& \ddots  & \vdots\\
		a_{n1}	& \dots   & a_{nn}\\
	\end{pmatrix}
	\\[0.4cm]
	\begin{bmatrix} 
		100&250\\
		300&499
	\end{bmatrix}
	\\[0.4cm]
	\begin{Bmatrix} 
		100&250\\
		300&499
	\end{Bmatrix}
	\\[0.4cm]
	\begin{Vmatrix}
		100&250\\
		300&499
	\end{Vmatrix}
\end{align}

\section{Klammern und Kästchen}

\begin{center}
	\( \Vert x\Vert_{p}=
	\left(
	\sum_{i=1}^{n} | x_{i} |^{p}
	\right)^{\frac{1}{p}} \)
\end{center}

\begin{center}
	\( \left.
	\begin{array}{lc|r}
		a&b&c\\
		\hline
		x&y&z\\
		c&a&b
	\end{array}
	\right\}
	\Rightarrow z,b \)
\end{center}

\begin{equation*}
\mbox{
	\boxed{
		\sin^2\varphi+\cos^2\varphi=1
	}
}
\end{equation*}


%\chapter{Programm- bzw. Quellcode}
\section{Quellcode}
Ein wichtiger Punkt ist auch, dass man Quellcode Stücke mit in seinen Praxisbericht einbaut. Hier nun einfach mal ein Beispiel in Form eines kleinen JAVA Codes, welcher aus einer Datei gelesen wird:

\lstinputlisting[
	label=code:algQuersumme,    % Label; genutzt für Referenzen auf dieses Code-Beispiel
	caption=Algorithmus zur Berechnung der Quersumme,
	captionpos=b,               % Position, an der die Caption angezeigt wird t(op) oder b(ottom)
	style=EigenerJavaStyle,     % Eigener Style der vor dem Dokument festgelegt wurde
	firstline=3,                % Zeilennummer im Dokument welche als erste angezeigt wird
	lastline=18                 % Letzte Zeile welche ins LaTeX Dokument übernommen wird
]{Quellcode/Eigenes-Java-File.java}

Zu beachten ist, dass jedes Stück Code kommentiert werden sollte. Was wird in diesem Abschnitt genau durchgeführt. Wo könnten Probleme auftreten und warum wurde dieses Stück hier hinzugefügt.

\vspace*{0.5cm}

\pagebreak
\textbf{EVENT 01} \textit{INIT}
\lstinputlisting[
	label={code:SSCUI-AfterInit},  % Label; genutzt für Referenzen auf dieses Code-Beispiel
	caption={Initialisierung im Programm},
	captionpos=b,               % Position, für die Caption:  t(op) oder b(ottom)
	style=EigenerABAPStyle,     % Eigener Style der vor dem Dokument festgelegt wurde
	firstline=4,                % Zeilennummer im Dokument welche als erste angezeigt wird
	lastline=23                 % Letzte Zeile welche ins LaTeX Dokument übernommen wird
]{Quellcode/Eigenes-ABAP-File.abap}

Es ist auch möglich, innerhalb des Listings \LaTeX{} Befehle zu verwenden. Dazu muss aber eine Escape-Sequenz angegeben werden. Das folgende Beispiel enthält ein Label, auf das dann verwiesen werden wird: Auch hier funktioniert \texttt{\textbackslash{}autoref}: \enquote{In \autoref{code:var_b} wird der Variable \texttt{b} \ldots}.

\begin{lstlisting}[
  style=EigenerJavaStyle,
  captionpos=b,
  caption={Zuweisung von Variablen},
  label={code:basic_block},
  escapeinside={@}{@}]
int a = 10;
@\label{code:var_b}@int b = a + 20;
return;
\end{lstlisting}

\pagebreak
Weitere Programmiersprachen können auch eingebunden werden. Hier mal ein Beispiel in der Programmiersprache Python:
\lstinputlisting[
	label=code:WhileLoop,    % Label; genutzt für Referenzen auf dieses Code-Beispiel
	caption=Algorithmus zum Schätzen einer Zahl in Python,
	captionpos=b,               % Position, an der die Caption angezeigt wird t(op) oder b(ottom)
	style=EigenerPythonStyle,   % Eigener Style der vor dem Dokument festgelegt wurde
	firstline=0,                % Zeilennummer im Dokument welche als erste angezeigt wird
	lastline=23                 % Letzte Zeile welche ins LaTeX Dokument übernommen wird
]{Quellcode/Eigenes-Python-File.py}

\clearpage % Absichtlicher Seitenumbruch, um ein besseres Layout zu erhalten.

\section{Pseudocode}
Pseudocode kann hilfreich sein, wenn \enquote{richtig} implementierte Algorithm in einer Programmiersprache zu lang sind und diese mittels Pseudocode bündig zusammengefasst werden können. \autoref{lst:euclid} zeigt Pseudocode.

Die Doku für das Package und damit eine Liste aller Befehle findet sich unter \newline
\url{http://tug.ctan.org/macros/latex/contrib/algorithmicx/algorithmicx.pdf}

\begin{algorithm}
	\caption{Euclid's algorithm}\label{lst:euclid}
	\begin{algorithmic}[1]
		\Procedure{Euclid}{$a,b$}\Comment{The g.c.d. of a and b}
			\State $r\gets a\bmod b$
			\While{$r\not=0$}\Comment{We have the answer if r is 0}
				\State $a\gets b$
				\State $b\gets r$
				\State $r\gets a\bmod b$
			\EndWhile\label{euclidendwhile}
			\State \textbf{return} $b$\Comment{The gcd is b}
		\EndProcedure
	\end{algorithmic}
\end{algorithm}

%\chapter{Literaturhinweise}

\textbf{Hinweis}: Verwendest du eine Quelle nicht, dann nimmt \LaTeX{} diese nicht mit ins Literaturverzeichnis auf!

\section{Fußnoten und Literaturverweise}
Hamburger, Döner, Currywurst\footnote{Hier fehlt eindeutig das Lieblingsessen der Informatiker, die Pizza!} - jeder kennt sie, jeder liebt sie und jeder isst sie.
Weil die Zeit drängt\cite{Bonnen.2016}, der Hunger groß ist und der nächste Schnellimbiss\cite{Forsthuber.2016} nur drei Schritte voraus. Und nach dem Essen?
Sind wir zwar satt, aber meist nicht wirklich glücklich, weil Fastfood\cite{Friedl.2017} meist eben auch nicht wirklich gut ist \cite{Kuhnlein.2016}.

Ja, uns ist bewusst, dass die Literatur nicht zum Text passt.
Deswegen hier nochmals der Rest.\cite{Mukherjee.2017, Preuss.2017, Schell.2017, Smith.2017, Visser.2017, Zaidi.2017, o.V..o.J.}

Übrigens: \acfp{ISBN} gehören nicht in die Bibliographie und werden deshalb auch nicht angezeigt.


% \include{Inhalt/04_Inhalt/einleitung.tex}
% \include{Inhalt/04_Inhalt/grundlagen.tex}
% \include{Inhalt/04_Inhalt/dynamischeLearningNFTs.tex}
% \include{Inhalt/04_Inhalt/anwendbarkeitfürXGP.tex}
% \include{Inhalt/04_Inhalt/schlussbetrachtung.tex}

\chapter{Einleitung}

abc
\chapter{Theoretische Grundlagen}

\section{Grundlagen der Betriebsplankosten- und -erfolgsrechnung (BPKER)}

Die BPKER, die das betriebliche Rechnungswesen mit den Planungs- und Überwachungsprozessen von Produktion und Absatz verknüpft, ist eine Weiterentwicklung der flexiblen Plankostenrechnung und Deckungsbeitragsrechnung. Hierbei sind Betriebs- und Absatzmodelle, die auf Modellen wie dem Hochofenprozessmodell basieren, entscheidend. Sie helfen, Einflussgrö{\ss}en auf den Periodenerfolg zu erfassen und berücksichtigen dabei nur die für die Planung, Überwachung und Produktkalkulation relevanten Beziehungen. Absatzmodelle, die auf Analysen vergangener Prozesse und Prognosen zukünftiger Absatzaktivitäten basieren, sind ebenso bedeutsam. Sie unterstützen marktsegmentspezifische Modellbildung und Prognosen zukünftiger Absatzentwicklungen. Des Weiteren ermöglichen sie fundierte Absatz- und Erlösplanungen, vorausgesetzt, man kennt die bedeutendsten Erlöseinflussgrö{\ss}en und ihre Wirkungszusammenhänge. \footcite[Vgl.][S. 295f]{Artikel_orginal}

Das primäre wirtschaftliche Ziel von Unternehmen ist die Maximierung des Periodenerfolgs. Die BPKER, die auf Betriebs- und Absatzmodellen basiert, zielt auf die Erreichung der optimalen Produktions- und Absatzergebnisse ab. Durch die Überwachung der Umsetzung dieser Pläne können Erwartungen für zukünftige Periodenkosten und -erlöse ermittelt werden. Während die flexible Plankosten- und Deckungsbeitragsrechnung nur die periodenbezogene Produktions- und Absatzplanung berücksichtigt, beinhaltet die BPKER zusätzlich die Planung und Überwachung der produktionsbezogenen Kosten und Absatzerlöse. Sie ermöglicht eine detailliertere Ermittlung der Ursachen von Plan-Ist-Abweichungen und erfüllt alle notwendigen Dokumentationsanforderungen. \footcite[Vgl.][]{Artikel_orginal}

\section{Aufbau der BPKER}

Die Erstellung einer BPKER beginnt mit der Identifizierung wesentlicher Kosten- und Erlöseinflussgrö{\ss}en, die die Ursachen für Kostenverbräuche oder Faktoreinsätze sind. Die Verfügbarkeit dieser Einflussgrö{\ss}en variiert, wobei einige intern kontrolliert und andere exogen bestimmt werden. Neben primären gibt es auch sekundäre Einflussgrö{\ss}en (Zwischenergebnisse in der betrieblichen Einflussgrö{\ss}enrechnung). Ihre Wirkungsweise auf Faktoreinsätze und Absatzleistungen kann durch statistische und analytische Verfahren bestimmt werden. Aufgrund der begrenzten Gültigkeit der statistisch ermittelten Koeffizienten, sind regelmä{\ss}ige Überprüfungen der Einflussgrö{\ss}enfunktionen nötig. Die Strukturmatrix hilft, Produktions- und Absatzprozesse abzubilden indem sie primäre Einflussgrö{\ss}en und abgeleitete Zielgrö{\ss}en darstellt. Sie verfolgt einen ähnlichen Aufbau, wobei das Absatzmengenprogramm im ersten Schritt bestimmt wird. Auftragsbestände können häufig als Grundlage für monatliche und quartalsbezogene Absatz- und Produktionsplanung dienen. Schlie{\ss}lich wird das Absatzprogramm in Bezug auf Vertriebswege, zusätzliche Dienstleistungen und erlösbestimmende Komponenten spezifiziert. \footcite[Vgl.][]{noauthor_betriebsplankostenrechnung_2023}

Betriebs- und Absatzmodelle, die als Input-Output-Modelle fungieren, stellen die Beziehungen zwischen Produktmengen und deren Einflussgrö{\ss}en dar. Mittels multiplikativer Verknüpfung mit Preisvektoren werden differenzierte Planerlöse und Plankosten je Periode generiert. Der Planperiodenerfolg entspricht der Differenz zwischen Gesamtperiodenerlösen und -kosten. Bei aufeinanderfolgenden Periodenmodellen werden Bestandszuführungen und Bestandsentnahmen von Halb- und Fertigfabrikaten berücksichtigt. Es besteht die Möglichkeit, sowohl Alternativkalkulationen durchzuführen als auch Optimierungsrechnungen mittels linearer Programmierungsansätze anzuwenden. \footcite[Vgl.][S. 302]{noauthor_betriebsplankostenrechnung_2023}

\section{Einsatz der BPKER im Unternehmenscontrolling bei Sorten- und Serienfertigung}

Betriebs- und Absatzmodelle reflektieren Mengenbeziehungen zwischen Produktionsprogramm und wesentlichen Einflussgrö{\ss}en. Änderungen im Produktionsprogramm erfordern Anpassungen wie Verfahrenswechsel oder Varianz von Fertigungsgrö{\ss}en. Dabei sind die Erfolgswirkungen verschiedener Ma{\ss}nahmen entscheidend. Die Mengenberechnung ermittelt Erzeugniseinsatzmengen und Planzeiten, und die Bewertungsrechnung berechnet Periodenkosten durch Anwendung geeigneter Preiskategorien auf Kostengüterbedarfsmengen. Die zunehmende Differenz zwischen Jahres-Plankosten und aktuellem Preisniveau erfordert kurz- und mittelfristige Anpassungen und Optimierungsansätze. Durch Simulationen können die Auswirkungen von Rabatten, Boni und Vertriebswegvariationen auf den Periodenerlös beurteilt werden. \footcite[Vgl.][S. 303ff]{Artikel_orginal}

Die Abweichungsanalyse ist ein Schlüsselwerkzeug für die kurzfristige Erfolgsüberwachung und kann mit unterschiedlicher Granularität durchgeführt werden. Sie hilft dabei, Ursachen und Verantwortlichkeiten für Abweichungen aufzudecken und gliedert Unterschiede zwischen Plan- und Isterfolg nach Erlös- und Kostenabweichungen. Änderungen in den Kostengüterabnahmen aufgrund von veränderten Umständen erfordern Anpassungen. Zudem sind Abweichungen von Sollgrö{\ss}en zu berücksichtigen, insbesondere wenn sie au{\ss}erhalb der statistischen Streubreite liegen. Am Monatsende können anhand der tatsächlichen Produktionsmenge und -bedingungen die Sollkosten berechnet und mit den Istkosten verglichen werden, was wesentliche Abweichungen verdeutlicht. Entscheidungsabhängige Abweichungen ergeben sich durch den Vergleich von Plankosten zu Monatsbeginn und den Sollkosten nach Monatsende. \footcite[Vgl.][S. 305ff]{Artikel_orginal}

Die Erzeugnis- und Auftragskosten und -erlöse sind wesentliche Bestandteile der Kostenträgerrechnung, die zur Preisbildung, Marktpreisbewertung, Lagerbewertung und Erfolgsmessung herangezogen werden. Die BPKER ermöglicht eine flexible Kostenabgrenzung der Kostenträgerrechnung, vorausgesetzt, man setzt für alle Freiheitsgrade einwertige Vorgaben ein. Alternativkalkulationen können in Betracht gezogen werden und Preisänderungen auf die Produktionskosten werden durch die trennende Mengen- und Bewertungsrechnung deutlich. In gleicher Weise erlauben Absatzmodelle Brutto- und Nettoumsatzberechnungen pro Verkaufseinheit. Durch Kombination von Erlös- und Kostenkalkulation wird eine Erfolgskalkulation pro Produkt oder Auftrag ermöglicht, wobei die periodenbezogene Produktions- und Absatzplanung keine Stückerfolgsberechnung erfordert. \footcite[Vgl.][S. 307f]{Artikel_orginal}

\section{Konzeptionelle Unterschiede zwischen BPKER und flexibler Plankostenrechnung}

Die flexible Plankosten- und Deckungsbeitragsrechnung, die auf der Proportionalität zwischen variablen Kosten und Produktmengenvariationen je Periode basiert, eignet sich vor allem für die Kalkulation, Planung und Kontrolle von Produktions- und Absatzprozessen. Alle anderen Kosten- und Erlöseinflussfaktoren werden im Vorfeld festgelegt. Die BPKER hingegen ermöglicht durch die Erfassung aller wesentlichen Kosten- und Erlösabhängigkeiten die Berechnung des Periodenerfolgs von Planungsvarianten und die detaillierte Untersuchung von Abweichungen, während die Mengen- und Bewertungsrechnung strikt getrennt und laufend mit aktuellen Preisen aktualisiert wird. Sie bietet den Vorteil, Abweichungen nach technischen Einzelfaktoren unter Verantwortlichkeit zu gliedern, während die flexible Plankostenrechnung eine globale Ursachenanalyse ermöglicht. \footcite[Vgl.][S. 308ff]{Artikel_orginal}

\section{Ergänzung der periodenbezogenen BPKR durch eine Online-Kennziffernrechnung}

In Zeiten steigender Automatisierung, hoher Anlagenintensität und sich ändernder Lohnformen nimmt auch der Anteil der Vorlaufkosten zu. Daher ist eine fein abgestimmte Planung der Kosten essentiell. Die BPKER begegnet diesem Bedarf, indem sie die Kosten- und Erlösabhängigkeiten eines Unternehmens umfassend erfasst und die Berechnung des Periodenerfolgs von Planungsvarianten ermöglicht. Die operative, zeitaktuelle Kennzahlenrechnung ermöglicht eine laufende Überwachung des Produktionsprozesses. Eine Online-Erfassung von Kosten, Erlösen und Mengen- undDasZeitgrö{\ss}en macht Abweichungen von Vorgabewerten sichtbar und ermöglicht rasche Anpassungen. k mit der üblichen Monatsrechnung verbundenes Online-Informationssystem unterstützt bei der Durchführung von Plan-/Soll- und Soll-/Istvergleichen. \footcite[Vgl.][S. 310ff]{Artikel_orginal}

\chapter{Praktischer Teil}


\section{Anwendung auf die HobelKunstwerkstatt GmbH}

\subsection{Unternehmensporträt}

Um die BPKER auf ein fiktives Unternehmensszenario anzuwenden, soll zuerst das betrachtete Unternehmen kurz skizziert werden. Die HobelKunstwerkstatt GmbH ist eine kleine Schreinerei mit zehn Mitarbeitern, die hochwertige Designermöbel herstellt. Das seit 90 Jahren in dritter Generation familiengeführte Unternehmen hat sich hauptsächlich auf zwei Produkte spezialisiert: Tische und Stühle aus Eichenholz. Da die Produkte im Premiumsegment angesiedelt sind, hat das Unternehmen im letzten Jahr einen Umsatz von 2,4 Mio. {\euro} erwirtschaftet. Hiervon entfällt der Gro{\ss}teil auf die beiden Hauptprodukte und ein kleiner Teil auf den Verkauf von Pflegeprodukten und die Aufbereitung von älteren Möbeln.

\subsection{Betriebs- und Absatzmodell}

\subsubsection{Betriebsmodell}

Ein Betriebsmodell definiert, wie ein Unternehmen seine Geschäftsprozesse und -strukturen für eine effiziente Wertschöpfung organisiert. Im Fall der HobelKunstwerkstatt GmbH sieht das Betriebsmodell folgenderma{\ss}en aus:

\begin{itemize}
    \item \textbf{Beschaffung:} Die HobelKunstwerkstatt arbeitet mit mehreren Lieferanten aus der Umgebung zusammen, um lokal nachhaltiges Holz zu beschaffen. Alle anderen Hilfsstoffe werden vom Gro{\ss}handel bezogen. 
    \item \textbf{Produktion:} Das Unternehmen produziert die beiden Hauptprodukte in der eigenen Produktionsstätte. Der Produktionsprozess besteht aus dem Verarbeiten des Rohmaterials in die Einzelteile der Produkte, den Zusammenbau der Einzelteile, dem Veredeln des Produkts und abschlie{\ss}end der Qualitätskontrolle.
    \item \textbf{Verwaltung/ Vertrieb:} Das Unternehmen hat eine Bürokraft, die sich zusammen mit dem Inhaber um sämtliche Verwaltungsaufgaben kümmert. Auf die Vertriebswege wird im Absatzmodell noch näher eingegangen.
\end{itemize}

\subsubsection{Absatzmodell}

Ein Absatz- oder Vertriebsmodell definiert wie ein Unternehmen seine Produkte an seine Kunden verkauft. Das Absatzmodell der HobelKunstwerkstatt GmbH wird im Folgenden skizziert:

\begin{itemize}
    \item \textbf{Direktverkauf:} Der Hauptumsatz des Unternehmens wird über den Direktverkauf erwirtschaftet. Das Unternehmen hat einen Showroom, in dem die angebotetenen Tische und Stühle in verschiedenen Varianten, sowie die Anpassungsmöglichkeiten für die Kunden präsentiert werden.
    \item \textbf{Online-Shop:} Zudem betriebt die HobelKunstwerkstatt einen Online-Shop, über den Kunden sich ihr Wunschprodukt digital konfigurieren und bestellen können. Zudem werden hierüber Pflegeprodukte vertrieben. Aufgrund der hochpreisigen Produkte ist der Online-Shop aber nur für einen kleinen Teil des Umsatzes verantwortlich.
\end{itemize}

\subsection{Einflussgrö{\ss}en auf das Betriebs- und Absatzmodell}

\subsubsection{Primäre Einflussgrö{\ss}en}

Primäre Einflussgrö{\ss}en sind Einflüsse, die einen direkten und wesentlichen Einfluss auf das Betriebs- und Absatzmodell haben. Änderungen dieser Einfussgrö{\ss}en resultieren meistens in erheblichen Schwankungen \zB der Produktkosten oder Verkaufserlöse. Deshalb ist es besonders wichtig diese konstant zu überwachen, um ggf. frühzeitig Anpassungen in den betroffenen Faktoren vornehmen zu können.

Im Betriebsmodell sind die Rohstoff- und Arbeitskosten die wichtigsten primären Einflussgrö{\ss}en. Diese machen zusammen beim Produkt Tisch 81,25 \% (Rohstoffe: 31,25 \% bzw. Arbeit: 50 \%) und beim Produkt Stuhl 83,33 \% (Rohstoffe: 33,33 \% bzw. Arbeit: 50 \%) aus. Das ist ein typisches Phänomen, da Deutschland ein Hochlohnland ist und somit der Arbeitslohn meist einen Gro{\ss}teil der Produktkosten ausmacht. Der hohe Einfluss der Rohstoffkosten ist hingegen eher durch den Fokus des Unternehmens, das verwendete Holz lokal und auch nachhaltigen Quellen zu beschaffen, zu erklären.

Im Absatzmodell sind der Preispunkt der Produkte und die Marktnachfrage die bedeutendsten primären Einflussgrö{\ss}en. Der Preispunkt hat einen erheblichen Einfluss auf die Verkaufszahlen und somit auf den Umsatz bzw. das Betriebsergebnis des Unternehmens. Gerade weil das Unternehmen Produkte im oberen Premiumsegment produziert, muss dieser sorgfältig gewählt werden, um die Balance zwischen Verkaufszahlen und Gewinnmarge zu halten. Auch die Marktnachfrage ist ein wichtiger Faktor, da die von dem Unternehmen verkauften Produkte eine sehr lange Lebenszeit haben und von den Kunden nicht nach einigen Jahren schon wieder erneuert werden. Jedoch ist dieser vom Unternehmen nicht komplett beeinflussbar, da die wirtschaftliche Situation au{\ss}erhalb seiner Kontrolle liegt und Kunden in schwierigen konjunkturellen Rahmenbedingungen Luxuskäufe eher noch zu Gunsten eines preiswerteren Produkts überdenken. 

\subsubsection{Sekundäre Einflussgrö{\ss}en}

Sekundäre Einflussgrö{\ss}en sind Faktoren, die das Betrieb- und Absatzmodell zwar beeinflussen, dies aber in eher indirekt und weniger stark. Zudem wirken sich diese Einflussgrö{\ss}en eher über Umwege bzw. in Kombination mit anderen Faktoren auf die Modelle aus. Somit sollten diese zwar auch regelmä{\ss}ig überwacht werden, jedoch nicht genauso wie die primären Einflussgrö{\ss}en im Zentrum stehen.

Im Betriebsmodell wären sekundäre Einflussgrö{\ss}en die Effizienz der Produktionsprozesse und die Qualität der Rohstoffe. Der erste Faktor beeinflusst, wie schnell und wirtschaftlich das Unternehmen produziert. Das wirkt sich somit einerseits intern auf die Zahlen des Unternehmens und andererseits extern auf Lieferzeiten und somit auch die Kundenzufriedenheit aus. Die Rohstoffqualität wirkt sich direkt auf die Qualität der Produkte aus und somit auch indirekt auf die Kundenzufriedenheit. Qualität ist gerade bei hochpreisigen Produkten ein wichtiger Faktor.

Als sekundäre Einflussgrö{\ss}en im Absatzmodell können das Marketing und der Kundenservice benannt werden. Marketing ist für die HobelKunstwerkstatt ein Hebel um die Bekanntheit der Marke und somit auch die Nachfrage nach den eigenen Produkten zu beeinflussen. Gerade im Premiumsegment ist auch der Kundenservice ein sehr wichtiger Faktor, da Kunden beim Kauf eines so teuren Produkts meist Dinge wie Lieferung und Montage und unkomplizierte Hilfe bei Problemen erwarten. Somit können auch Weiterempfehlungs- und Wiederverkaufsraten gesteigert werden.

\subsection{Einflussgrö{\ss}enfunktionen}

Im Folgenden sollen noch die Auswirkungen der Einflussgrö{\ss}en auf verschiedene Aspekte des Betriebs- und Absatzmodells in einer Strukturmatrix veranschaulicht werden.

Legende:
\begin{itemize}
    \item \textbf{+:} positive Beziehung
    \item \textbf{-:} negative Beziehung
    \item \textbf{0:} keine Beziehung
\end{itemize}

\subsubsection{Betriebsmodell}

\begin{table}[H]
    \centering
    \begin{tabular}{|c|c|c|c|}
      \hline
      \textbf{Einflussgrö{\ss}e} & \textbf{Produktionskosten} & \textbf{Produktqualität} & \textbf{Produktionskapazität} \\
      \hline
      \underline{Primär} & & & \\
      \hline
      Rohstoffkosten & + & + & - \\
      \hline
      Arbeitskosten & + & + & - \\
      \hline
      \underline{Sekundär} & & & \\
      \hline
      Prozesseffizienz & - & + & + \\
      \hline
      Rohstoffqualität & + & + & - \\
      \hline
    \end{tabular}
    \caption{Strukturmatrix für das Betriebsmodell}
  \end{table} 

\subsubsection{Absatzmodell}

\begin{table}[H]
    \centering
    \begin{tabular}{|c|c|c|c|}
      \hline
      \textbf{Einflussgrö{\ss}e} & \textbf{Verkaufszahlen} & \textbf{Umsatz} & \textbf{Kondenzufriedenheit} \\
      \hline
      \underline{Primär} & & & \\
      \hline
      Preispunkt & - & + & - \\
      \hline
      Marktnachfrage & + & + & 0 \\
      \hline
      \underline{Sekundär} & & & \\
      \hline
      Marketing & + & + & + \\
      \hline
      Kundenservice & + & + & + \\
      \hline
    \end{tabular}
    \caption{Strukturmatrix für das Absatzmodell}
  \end{table} 

\subsection{Betriebsplankosten- und -erfolgsrechnung}

Die Kostenstruktur für die Herstellung dieser Produkte setzt sich vereinfachend aus Materialkosten, Arbeitskosten und dem Deckungsbeitrag für alle sonstigen Kosten zusammen. In der folgenden Tabelle werden diese Kosten genauer aufgeschlüsselt:

\begin{table}[H]
    \centering
    \label{tab:Kostenstruktur_Produkte}
    \begin{tabular}{|c|c|}
      \hline
      Tisch & Stuhl \\
      \hline
      Arbeitskosten: 4.000 {\euro} & Arbeitskosten: 600 {\euro} \\
      Materialkosten: 2.500 {\euro} & Materialkosten: 400 {\euro} \\
      Deckungsbeitrag: 1.500 {\euro} & Deckungsbeitrag: 200 {\euro} \\ 
      \hline
      Kosten insgesamt: 8000 {\euro} & Kosten insgesamt: 1.200 {\euro} \\
      \hline
      Nettoverkaufspreis: 9.600 {\euro} & Nettoverkaufspreis: 1.440 {\euro} \\
      \hline
      Bruttoverkaufspreis: 11.424 {\euro} & Bruttoverkaufspreis: 1713,60 {\euro} \\
      \hline
    \end{tabular}
    \caption{Kostenstruktur der Produkte}
  \end{table}

Für die konkrete Anwendung soll jetzt angenommen werden, dass das Unternehmen im Monat zehn Tische und 40 Stühle fertigt. Da die Gewinnmarge bei beiden Produkten 20 \% und die Kunden aufgrund der abgestimmten Optik meist einen Tisch zusammen mit vier Stühlen kaufen, würde es auch nicht sinnvoll sein, die vorhandenen Produktionskapazitäten zugunsten des Produkts Tisch zu verschieben, auch wenn bei diesem der Gewinn absolut höher ausfallen würde.

Mit diesen Informationen kann die BPKR verwendet werden, um die Gesamtkosten für die Produktion im nächsten Monat zu berechnen.:

\[ Kosten Tische: 8.000 EUR * 10 = 80.000 EUR \]
\[ Kosten Stühle: 1.200 EUR * 40 = 48.000 EUR \]
\[ Gesamtkosten: 80.000 EUR + 48.000 EUR = 128.000 EUR \]

Die BPKR bietet jedoch nicht nur eine Momentaufnahme der Kostenstruktur, sondern auch eine Grundlage für die BPER. Hierfür sollen im Folgenden zuerst die Gesamterlöse des nächsten Monats berechnet werden.

\[ Erlöse Tische: 9.600 EUR * 10 = 96.000 EUR \]
\[ Erlöse Stühle: 1.440 EUR * 40 = 57.600 EUR \]
\[ Gesamterlöse: 96.000 EUR + 57.600 EUR = 153.600 EUR \]

Mit den Ergebnissen BPKR lässt sich nun mittels der BPER das Betriebsergebnis durch Abziehen der Kosten von den Erlösen ermitteln:

\[ 153.600 EUR - 128.000 EUR = 25.600 EUR \]

Somit würde das Unternehmen im nächsten Monat einen Gewinn von 25.600 {\euro} erzielen, wenn alle Produktions- und Verlaufspläne genauso umgesetzt werden.

Basierend auf diesen Daten lässt sich nun auch ohne weiteres eine Abweichungsanalyse machen, falls das Betriebsergebnis nicht erreicht wurde, da ermittelt werden kann, durch welche Komponenten des Betriebs- oder Absatzmodells die Abweichungen verursacht wurden und dann mittels der Strukturmatrix die Einflussgrößen, die verändert werden müssen, bestimmt werden können. Wenn diese dann bekannt sind, können konkrete Maßnahmen beschlossen werden, um das Betriebsergebnis wieder in die gewünschte Richtung zu korrigieren.
\chapter{Fazit/ Schlussbetrachtung}

abc

% ---- Literaturverzeichnis
\cleardoublepage
\renewcommand*{\chapterpagestyle}{plain}
\pagestyle{plain}
%\pagenumbering{Roman}                   % Römische Seitenzahlen
%\setcounter{page}{\numexpr\value{savepage}+1}

\printbibliography[title = Literaturverzeichnis]

% ---- Anhang
\appendix
%\clearpage
%\pagenumbering{Roman}  % römische Seitenzahlen für Anhang

%\chapter{Anhang}

% \section{Fragebögen Experteninterviews}

% \subsection{Workflows} % \label{FB_Workflows}

abc

%\includepdf[pages = -]{Literatur/bgPF_Wiki.pdf}

\newpage
\end{document}
